% Generated draft: requires mathematical and source review.
\documentclass[11pt]{article}
\usepackage[margin=25mm]{geometry}
\usepackage{fontspec}
\setmainfont{Latin Modern Roman}
\usepackage{amsmath,amssymb,mathtools,unicode-math}
\setmathfont{Latin Modern Math}
\usepackage{enumitem,needspace,xurl,hyperref}
\setlength{\parindent}{0pt}
\setlength{\parskip}{5pt}
\setlength{\emergencystretch}{3em}
\newcommand{\N}{\mathbb{N}}
\newcommand{\Z}{\mathbb{Z}}
\newcommand{\Q}{\mathbb{Q}}
\newcommand{\R}{\mathbb{R}}
\newcommand{\C}{\mathbb{C}}
\newcommand{\F}{\mathbb{F}}
\newcommand{\A}{\mathbb{A}}
\newcommand{\PP}{\mathbb{P}}
\newcommand{\E}{\mathbb{E}}
\newcommand{\eps}{\varepsilon}
\newcommand{\dd}{\,\mathrm{d}}
\begin{document}
\section*{UnsolvedMath ID: 701582}
\textbf{Source problem code:} catalogue entry
\subsection{Definitions and mathematical statement}
Let \(G\) be a group, and use the commutator convention \([x,g]=x^{-1}g^{-1}xg\). For each \(g\in G\), define the set of commutator values
\[
C_g=\{[x,g]:x\in G\}.
\]
The hypothesis concerns this set itself being a subgroup, not merely the subgroup it generates. Define the lower central series by
\[
\gamma_1(G)=G,\qquad
\gamma_{n+1}(G)=\langle [u,v]:u\in\gamma_n(G),\ v\in G\rangle
\quad(n\geq 1).
\]
The group \(G\) is residually nilpotent if
\[
\bigcap_{n\geq 1}\gamma_n(G)=\{1\}.
\]
Equivalently, every nonidentity element has nonidentity image under some homomorphism from \(G\) to a nilpotent group; those target groups need not be finite.

\emph{Question (Kourovka Notebook, Problem 19.34).} Is it true that every finitely generated group \(G\) satisfying
\[
(\forall g\in G)\quad C_g\leq G
\]
is residually nilpotent?

The supplied record additionally reports that the answer is affirmative when \(G\) is finite. This finite-case assertion is retained as a source-reported qualification and has not been independently verified. The record labels the general question open; this is supplied registry status, not an independently verified literature assessment. The attribution is malformed in the supplied record and remains unresolved.
\subsection{Short English statement}
Must a finitely generated group be residually nilpotent if, for each fixed element \(g\), all commutators \([x,g]\) form a subgroup as \(x\) varies? The supplied record reports an affirmative answer for finite groups and labels the general question open.
\subsection{Sources}
\begin{itemize}
\item Local UnsolvedMath registry, permanent ID \(701582\), exact source code \texttt{catalogue entry}, title \emph{Kourovka Notebook Problem 19.34}. Supplied upstream revision: \texttt{d371ee917ac9647fbb06cd6fc529deec6f3dda82}. The record was updated on \texttt{2026-10-16T14:00:00Z}; this timestamp is registry metadata, not evidence of independent literature verification.
\item The supplied background identifies \emph{The Kourovka Notebook (Unsolved Problems in Group Theory)}, 19th issue (2018), original Problem 19.34, and the 21st issue edition (2026), identified as \texttt{arXiv:1401.0300v48}. These bibliographic details are supplied source attribution and have not been independently checked. The collection description names E. I. Khukhro and V. D. Mazurov as editors.
\item Attribution is corrupted: the supplied proposer field contains the literal fragment \texttt{D. Gon\textbackslash{}c\{c}; the background repeats that fragment and places \texttt{alves, T. Nasybullov} beneath a discussion-and-literature label. The intended complete attribution cannot be recovered unambiguously from the supplied fields.
\item The supplied statement reports the affirmative finite case without a separate supporting citation. No source was visited for this draft.
\end{itemize}
\end{document}
